\chapter{Discussão e limitações conhecidas}
\label{cap:discussao}

Este capítulo consolida, de forma explícita, as simplificações
deliberadas e limitações conhecidas da implementação atual -- distintas
de bugs (já corrigidos ao longo dos capítulos anteriores), no sentido de
que são escolhas de escopo documentadas, com risco avaliado e aceito
para o caso de estudo validado, mas que podem se tornar relevantes em
outros contextos de uso.

\section{Mistura de terreno de fronteira não implementada}
\label{sec:blend_bdy_terrain}

A limitação de maior impacto numérico mensurável identificada é a
ausência do passo \texttt{config\_blend\_bdy\_terrain}: no algoritmo de
referência, antes de gerar a grade vertical (\cref{cap:fase2}), a altura
de terreno das células do anel de fronteira da malha regional
(\texttt{bdyMaskCell}~$\ge 1$) é misturada com o terreno do dado de
\textit{fundo} (\textit{first-guess}), suavizando a transição entre o
domínio regional e a condição de contorno que o alimentará. Sem esse
passo, o resíduo de $\sim\SI{171}{\metre}$ documentado na
\cref{sec:validacao_fase2} permanece concentrado exatamente nessa faixa
de células -- que é, por definição, onde a fronteira lateral atua. Para
o caso de validação usado neste relatório, esse resíduo não impediu uma
integração de 24 horas fisicamente sã (\cref{cap:resultados}), mas é a
lacuna mais provável de precisar de atenção em usos futuros que dependam
de maior precisão perto da borda do domínio regional, ou em domínios com
relevo mais acentuado próximo à fronteira do que o caso testado.

\section{Reamostragem vertical de solo simplificada}

O algoritmo de referência reamostra os perfis de temperatura e umidade
de solo do \textit{first-guess} para as quatro profundidades-padrão do
esquema Noah usando uma interpolação linear por profundidade própria
(distinta e independente da grade vertical atmosférica). Esta
implementação copia esses campos diretamente do \textit{first-guess}
sem essa reamostragem -- uma simplificação segura apenas porque, no caso
de uso deste projeto, tanto a origem quanto o destino do
\textit{first-guess} são o próprio \mpasa{}, com o mesmo esquema Noah e
as mesmas quatro profundidades de solo. Essa simplificação deixa de ser
válida se o \textit{first-guess} vier de uma fonte com profundidades de
solo diferentes (por exemplo, GFS via WPS/\texttt{ungrib}).

\section{Reclassificação de gelo marinho não implementada}

O algoritmo de referência reclassifica células de água muito fria
(temperatura de pele abaixo de um limiar configurável,
\texttt{config\_tsk\_seaice\_threshold}) como gelo marinho ou terra,
ajustando em cascata o tipo de uso do solo, textura de solo, albedo
máximo de neve, e os perfis de temperatura/umidade de solo dessas
células. Essa reclassificação não foi implementada. A omissão é
irrelevante para o caso de validação (América do Sul: o campo
\texttt{SEAICE} do \textit{first-guess} é identicamente zero em toda a
malha, confirmado diretamente contra o dado real), mas se tornaria
relevante em qualquer malha regional que inclua regiões de alta
latitude com gelo marinho real.

\section{Umidade específica via umidade relativa não implementada}

O código de referência oferece dois caminhos para diagnosticar a razão
de mistura de vapor d'água: a partir da umidade específica do
\textit{first-guess} (\texttt{config\_use\_spechumd=.true.}) ou, na
ausência desse campo, a partir da umidade relativa via a função de
saturação de Flatau/Thompson (\texttt{rslf}). Apenas o primeiro caminho
foi implementado -- suficiente porque todos os casos de estudo deste
projeto, até o momento, usam \texttt{config\_use\_spechumd=.true.} com um
\textit{first-guess} que de fato contém esse campo. O programa detecta
explicitamente a configuração oposta e encerra com uma mensagem de erro
clara, em vez de produzir um resultado silenciosamente incorreto.

\section{Escopo de validação}

Todo o processo de validação numérica e funcional apresentado neste
relatório usa uma única malha regional (\texttt{SouthAmerica},
$\sim\SI{60}{\kilo\metre}$), um único caso de estudo (24 horas de
integração a partir de \texttt{2026-01-01\_00:00:00}), e
$n_{VertLevels}=55$. A lista de 61 níveis de pressão fixa usada como
representação intermediária entre a extração de campos e a interpolação
horizontal nativa (\texttt{pressure\_levels.F90}, herdada sem
modificação da rota original) é derivada de medianas estatísticas de
altura de nível de modelo -- não específica da América do Sul, e portanto
esperada como razoável em outras malhas -- mas nunca de fato testada em
uma malha/resolução distinta. A generalização estrutural do código
(\cref{sec:generalizacao}) remove a maior barreira à reprodução em outros
casos, mas não substitui uma validação explícita nesse cenário mais
amplo.

\section{Um interpolante linear por partes: limite de suavidade}

A interpolação baricêntrica adotada (\cref{sec:baricentrica}) é
$C^0$ -- contínua entre triângulos adjacentes do dual de Delaunay, mas
com derivada descontínua nas arestas desses triângulos. Para o problema
resolvido aqui (gerar uma condição inicial/de fronteira, consumida
integralmente pelo núcleo dinâmico do modelo, que já aplica sua própria
suavização/filtro numérico durante a integração), essa limitação não se
manifestou como problema na validação funcional (\cref{cap:resultados}).
Ela poderia se tornar relevante caso o mesmo motor de interpolação fosse
reaproveitado para outro uso que dependa diretamente de gradientes
horizontais suaves do campo interpolado -- caso em que os interpolantes
de maior continuidade discutidos na revisão de literatura
(\cref{cap:literatura}, Hiyoshi e Sugihara~\cite{hiyoshi2000continuity})
seriam um candidato natural de extensão.

\section{Síntese das limitações}

\begin{table}[H]
\centering
\caption{Síntese das simplificações/limitações conhecidas, seu impacto
  no caso validado, e a condição sob a qual se tornariam relevantes.}
\label{tab:limitacoes}
\small
\begin{tabular}{@{}p{4.3cm}p{3.6cm}p{5.2cm}@{}}
\toprule
Limitação & Impacto no caso validado & Torna-se relevante quando \\
\midrule
Sem mistura de terreno de fronteira & Resíduo de $\sim\SI{171}{m}$ concentrado no anel de fronteira & Precisão perto da borda regional importa; relevo acentuado na fronteira \\
Solo copiado sem reamostragem vertical & Nenhum (mesma malha/esquema na origem e destino) & \textit{First-guess} com esquema de solo/profundidades distintas \\
Sem reclassificação de gelo marinho & Nenhum (sem gelo marinho no caso) & Malha regional em alta latitude \\
Sem \texttt{qv} via RH & Nenhum (\texttt{config\_use\_spechumd} sempre \texttt{.true.} nos casos testados) & \textit{First-guess} sem umidade específica direta \\
Validação em 1 malha/1 caso & -- & Uso em outra resolução/domínio/estação do ano \\
Interpolação $C^0$ (baricêntrica linear) & Nenhum observado & Uso que dependa de gradiente horizontal suave do campo interpolado \\
\bottomrule
\end{tabular}
\end{table}
